\begin{trapbox}
	\begin{description}[style=nextline, leftmargin=2.8em, labelsep=0.5em, itemsep=0.35em]
		\item[\textbf{\textcolor{CTrap}{易错点一（忽略分母为零）}}] 学生经常直接使用结论$\frac{S_{2m}-S_m}{S_m}=q^m$。但忘记条件$S_m\neq 0$。从而在$S_m=0$时错误应用。
		\item[\textbf{\textcolor{CTrap}{正确理解（检查分母）：}}] 应用性质前必须验证$S_m\neq 0$。若$S_m=0$，则分式无意义，结论不直接成立。
		\item[\textbf{\textcolor{CTrap}{错因分析（忽视前提）：}}] 学生只机械记忆公式，忽略条件，导致解题时未排除$S_m=0$的情况。
	\end{description}

	\medskip
	\begin{description}[style=nextline, leftmargin=2.8em, labelsep=0.5em, itemsep=0.35em]
		\item[\textbf{\textcolor{CTrap}{易错点二（错用片段和）}}] 误认为$S_m,\ S_{2m},\ S_{3m}$成等比。而实际是$S_m,\ S_{2m}-S_m,\ S_{3m}-S_{2m}$成等比。
		\item[\textbf{\textcolor{CTrap}{正确理解（明确对象）：}}] 性质针对的是每隔$m$项的“片段差”，即$S_{2m}-S_m$等，而不是前$n$项和$S_n$本身。
		\item[\textbf{\textcolor{CTrap}{错因分析（概念混淆）：}}] 学生未理解“片段和”的具体定义，混淆$S_{2m}$与$S_{2m}-S_m$。
	\end{description}

	\medskip
	\begin{description}[style=nextline, leftmargin=2.8em, labelsep=0.5em, itemsep=0.35em]
		\item[\textbf{\textcolor{CTrap}{易错点三（未讨论公比为一）}}] 条件中要求$q\neq 1$。但问题中出现了$q=1$的情况。学生直接套用求和公式$S_n=\frac{a_1(1-q^n)}{1-q}$。导致分母为零。
		\item[\textbf{\textcolor{CTrap}{正确理解（分类讨论）：}}] 当$q=1$时，数列为常数列$a_1$。前$n$项和为$S_n=na_1$。此时片段和性质仍然成立，但需单独推导，不能直接使用结论中的公式。
		\item[\textbf{\textcolor{CTrap}{错因分析（条件遗漏）：}}] 学生忽略了结论的适用条件，对$q$的取值不加讨论，或盲目套用导致错误。
	\end{description}
\end{trapbox}
